\section{Discussion}

\subsection{Formulation Choices and Conservative Structure}

The central modelling choice in this work is to use one fixed porosity field,
$\theta$, in the finite-volume transport equation and in the PHREEQC cell
volumes alike, which keeps aqueous accumulation, transport, and equilibrium
chemistry on the same volumetric basis. The reactive surface is prescribed
independently through the explicit cell field $a_{\mathrm{surf}}$, so the
chemical storage volume and the amount of reactive surface remain separately
defined material properties.

This approach extends earlier micro-continuum reactive transport formulations
to the case where the reactive surface area is an unresolved cell property
rather than an explicitly meshed boundary. In existing pore-scale DBS solvers,
including \nolinkurl{porousMedia4Foam} and earlier GeoChemFoam solvers,
reactive uptake is typically imposed as a flux condition on the diffuse
fluid-solid interface. This requires the reactive surface to coincide with the resolved image voxel
boundary, which is a good approximation when the dominant reactive area lies at
the contact between macropore and grain. In multi-scale carbonate images the
microporous phase holds a large internal reactive area that lies inside the
coarse cell rather than on its boundary. The cell-based formulation carries
that area as an intrinsic volumetric density
$a_{\mathrm{surf}}^{\mathrm{intrinsic}}$, scaled by the local phase fraction
and passed to PHREEQC as a cell-level property.

Conservation is therefore enforced at the cell level rather than through a face
flux. The combined inventory $M = \theta C + a_{\mathrm{surf}}\Gamma$ is updated
locally by PHREEQC at each timestep, and the aqueous transport step is
conservative in the standard finite-volume sense, so the only mass-balance
error is the splitting error of the coupling itself. The 1D analytical
benchmark measures that error directly.

\subsection{Geometric Scalability}

Because the validation ladder separates conservative coupling, surface-area
geometry, and microporosity, each error can be attributed to the layer that
produced it: the coupling to within $1.2\%$ in bulk inventory against the
analytical retardation factor, the interface-area construction to sub-percent
accuracy against analytic spheres and posts, and the intrinsic micropore
density only in the cases that use it.

The transition from 2D to 3D introduces qualitatively different transport paths
and surface-loading distributions. In the 2D case, open fractures provide
preferential pathways for dissolved Pb$^{2+}$, while strong-site Pb occupancy
varies across the surrounding microporous matrix and is enhanced near the inlet
and fracture intersections. In the 3D cases, the tortuosity of the
pore network produces a more dispersed front and a more heterogeneous
surface-loading pattern. This spatial heterogeneity in surface uptake controls the effective retardation
factor seen at the continuum scale, and following it requires a simulation that
resolves the local connectivity of the image.

The direct resolved-to-DBS comparison shows that the coarse representation also
preserves the bulk transport and sorption response across $Pe = 0.1$ to $10$,
with inventory differences below $0.9\%$ and profile errors below $5.8\%$, both
growing with $Pe$ across the sweep. That growth sets the direction of the
remaining uncertainty: higher $Pe$, different connectivity, or sharper kinetic
fronts may expose an upscaling failure that this range does not.

\subsection{Computational Cost and Scaling}

PHREEQC is called once per active cell per timestep, with a full speciation
and surface-equilibration calculation for each cell. The chemistry overhead
therefore depends primarily on the number of active cells, the geochemical
system, and the frequency of the chemistry call. For larger grids or more complex geochemical systems, masking near-solid cells
out of the PhreeqcRM mapping (here at
$\theta \le \theta_{\mathrm{mask}} = 10^{-3}$) is the main lever on that cost.
The threshold is held fixed across all examples here, and its influence on
accuracy and cost deserves a dedicated sensitivity study.

\subsection{Limitations and Future Work}

The intrinsic micropore surface density $a_{\mathrm{surf}}^{\mathrm{intrinsic}}$
is a user input per image phase rather than a quantity inferred from the
segmented image, so it must be estimated from complementary measurements such
as nitrogen BET adsorption, mercury intrusion porosimetry, or secondary image
analysis at a finer resolution \citep{noiriel2009changes}. Once the
surface-area construction has been validated, this density is the leading
physical uncertainty in microporous applications, because the accuracy of the
predicted uptake follows directly from how well it represents the internal
reactive area. Linking image-derived phase fractions to measured or computed
internal surface-area densities, and quantifying the sensitivity to that input,
is the next step.

The upscaled DBS treatment assumes that unresolved matrix properties can be
carried as local cell-averaged quantities. This holds when diffusive exchange
between the resolved pore space and the unresolved matrix is fast compared with
the advective residence time, so that the intrinsic surface area acts as a
locally accessible reactive reservoir. At larger $Pe$, advective bypass and
finite matrix-exchange times leave only part of that area accessible to a
passing concentration front, and the sorption error then carries uncertainty in
the temporal accessibility of the surface as well as in its amount. Local
equilibrium upscaling of an unresolved reactive matrix shares this limitation,
so $Pe$-resolved benchmark comparisons are needed when the method is applied to
strongly advective systems.

Transport and chemistry are coupled by a first-order Lie splitting, which
introduces a splitting error that scales with the timestep and the stiffness of
the reaction terms. For the surface-complexation systems used here, that error
sits within the benchmark agreement reported in Section~\ref{sec:validation},
but it may become significant for faster reactions or stiffer geochemical
systems. A second-order Strang splitting \citep{strang1968construction} is the
natural extension.

Restart depends on the serialised PHREEQC state written with each complete
checkpoint, because the finite-volume \texttt{Surf\_*} fields alone cannot
reconstruct that state. The production $Pe$ sweep restarted from these
checkpoints across several jobs, and moving a run between machines requires the
field data and its matching PHREEQC dump to travel together.

The cases here use a fixed surface area and equilibrium surface complexation,
and cover competitive multi-component chemistry and a Li-selective multi-ion
brine within that scope. Evolving reactive area, mineral dissolution and
precipitation, kinetic surface reactions, and the wider range of mineral
surfaces and aqueous compositions found in natural systems lie outside it.
Extending the model to them will require benchmarks that separate transport
error, surface-area evolution, and reaction-model uncertainty while retaining
cell-level mass-balance checks.
